% =============================================================================
\section{Concepts-Binding Coupling}
\label{sec:concepts}
% =============================================================================
Generic programming identifies an abstraction that consists of a
formal hierarchy of polymorphic abstract requirements on data types,
referred to as \emph{concepts}, and a set of classes that conform
precisely to the specified requirements, referred to as
\emph{models}~\cite{a-gps-99}. A hierarchy of related concepts can be
viewed as a directed acyclic graph, where a node of the graph
represents a concept and an arc represents a refinement relation. An
arc directed from concept A to concept B indicates that concept B
refines concept A. When a class template is instantiated, each one of
its template parameters is substituted with a model of one or more
concepts associated with the template parameter.

As mentioned in the introduction, \cgal{} rigorously adheres to the
generic programming paradigm. As a consequence, most components of
\cgal{} are either class or function templates. Many of the parameters
of these templates are described in terms of models of concepts. When
a class or a function template is instantiated, each one of its
template parameters is substituted with a model of one or more
concepts associated with the template parameter. Close to~750 concepts
can be identified in \cgal{} at the time this article is written. Most
hierarchy graphs of concepts are small. Few graphs, such as the graph
of concepts of the geometry traits of the \iiDArrangementsPackage{}
package is quite large with intricate refinement relations. Therefore,
these concepts are organized in clusters that describe the refinement
relations among the concepts within a cluster. We have introduced
tight coupling between concepts and (i) binding generations and (ii)
type annotation. We describe these relations in this section.

% -----------------------------------------------------------------------------
\subsection{Following Generic Concepts}
\label{sec:code:concepts}
% -----------------------------------------------------------------------------
For each concept we have introduced a template function that generates
bindings for all the type and function members required by the
concept. We call this function for every model of the concept. This
systematic approach guaranties that all the documented functions and
types, but nothing else, are exposed. Consider for example the cluster
of geometry traits concepts depicted in
Figure~\ref{fig:atc:central}. The template functions
\cppLstinline{export_AosBasicTraits\_2()},
\cppLstinline{export_AosXMonotoneTraits\_2()}, and
\cppLstinline{export_AosTraits\_2()} accept a binding class
object\footnote{A binding class object of type
  \cppLstinline{bp:class_<T>} is used to expose the C++ type
  \cppLstinline{T} to Python.} for a particular geometry traits model
and populate it with all attributes that correspond to the
requirements of the concepts \cppLstinline{AosBasicTraits\_2},
\cppLstinline{AosXMonotoneTraits\_2}, and \cppLstinline{AosTraits\_2},
respectively. If a concept \cppLstinline{B} refines a concept
\cppLstinline{A}, then the function \cppLstinline{export_B()} calls
the function \cppLstinline{export_A()}. The bindings of different
traits models are generated in separate compilation units. As nodes in
the concept refinement graphs may have more than a single parent, we
ensure that a function that corresponds to a concept is not invoked
more than once in a given compilation unit; see Lines~3--4 in the
definition of the function \cppLstinline{export_AosBasicTraits\_2()}
in Listing~\ref{lst:cpp:exportBasic}.

\begin{lstfloat}
  \begin{lstlisting}[style=cppFramedStyle,basicstyle=\normalsize,%
                     label={lst:cpp:exportBasic},%
                     caption={Exposing attributes that correspond to the concept \cppLstinline{AosBasicTraits_2}.}]
template <typename GeometryTraits_2, typename ClassObject, typename Concepts>
void export_AosBasicTraits_2(ClassObject c, Concepts& concepts) {
  // Sentinel
  static bool exported = false;
  if (exported) return;

  // Expose attributes that correspond to the traits
  auto& classes = concepts.m_basic_traits_classes;
  classes.m_point_2 = bp::class_<Point_2>("Point_2")
    .def(bp::init<>())
    .def(bp::init<const Point_2&>());
  ~~$\cdots$~~
}
  \end{lstlisting}
\end{lstfloat}

Typically, a concept in our case requires the provision of nested
types that are themselves models of other concepts. For example, the
concept \cppLstinline{AosBasicTraits\_2} requires the provision of the
nested type \cppLstinline{Point_2} used to represent a two-dimensional
point. The type \cppLstinline{Point_2} must be
\cppLstinline{DefaultConstructible} and
\cppLstinline{CopyConstructible}. Let \cppLstinline{p} be an object of
type \cppLstinline{Point_2} required by the
\cppLstinline{AosBasicTraits\_2} concept. Most traits models also
support the calls \cppLstinline{p.x()} and \cppLstinline{p.y()}; they
return the $x$ and $y$ exact Cartesian coordinates of the points $p$,
respectively. However, the \cppLstinline{Point_2} type nested in the
traits model \cppLstinline{Arr_Bezier_curve_traits\_2} does not
maintain an exact representation of the coordinates. Thus, it instead
supports the call \cppLstinline{p.approximate()}, which returns a pair
of floating point numbers that only approximate the real
coordinates. The function \cppLstinline{export_AosBasicTraits\_2()}
constructs a class object for the \cppLstinline{Point\_2} type, but
adds nothing but Python attributes for the default and copy
constructors of this type; see Lines~9--10 in the definition of the
function \cppLstinline{export_AosBasicTraits\_2()} above. It also
saves the class object of the \cppLstinline{Point_2} type (see
Line~7), so that later on the functions responsible for the bindings
of the various traits models can add additional Python attributes
(such as the \cppLstinline{x()}, \cppLstinline{y()}, and
\cppLstinline{approximate()} functions) to the class object.

The concept \cppLstinline{AosHorizontalSideTraits\_2} requires the
provision of the \cppLstinline{Compare\_ x\_on\_boundary\_2}
functor,\footnote{A functor is a class that has one or more
  \cppLstinline{operator()} members; thus, it acts as a function.}
which has two versions of the overloaded members
\cppLstinline{operator()} as follows; one accepts two ends of curves
(refer to the manual for the precise encoding of a curve end) and one
accepts an end of a curve end and a point. The former compares the
$x$-coordinates of the curves at their respective limits or ends and
the latter compares the $x$-coordinates of the point and the
$x$-coordinates of the curve at its respective limit or end.
%
The concept \cppLstinline{AosIdentifiedHorizontalTraits\_2} indirectly
refines the concept \cppLstinline{AosHorizontalSideTraits\_2}
mentioned above. (The former would be used to define a traits model
that handles arcs on a sphere with poles on the left and right sides;
Figure~\ref{fig:atc:sphericalBoundary} depicts a cluster of concepts
for models that handle arcs on a sphere with poles on the top and
bottom sides.) In addition to the two versions of the overloaded
members \cppLstinline{operator()} of the functor
\cppLstinline{Compare\_x\_on\_boundary\_2} it requires the provision
of a third version that accepts two points and compares their $x$
coordinates. Similar to the case described in the previous paragraph,
where the class object that handles the bindings for the
\cppLstinline{Point_2} type must be available in two separate
functions, also here the class object that handles the bindings for
the functor \cppLstinline{Compare\_x\_on\_boundary\_2} must be
available in both functions
\cppLstinline{export_AosHorizontalSideTraits\_2()} and
\cppLstinline{export_AosIdentifiedHorizontalTraits\_2()}. Saving the
class object in one function so that it can be reused latter on in the
other addresses this necessity.

% -----------------------------------------------------------------------------
\subsection{Type Annotation}
\label{ssec:concepts:annotation}
% -----------------------------------------------------------------------------
\cgal{} is a large library to start with. It is divided into
approximately~150 packages with countless function and class
templates. The vast number of function and type instances that can be
defined using this templates can be overwhelming. The development of
non critical code that is based on \cgal{} can be expedited using
Python and \cgal{} Python bindings.  Using code completion (which is
based on type annotations) offered by several Python IDEs can
accelerate the process much further. Type annotation in Python refers
to annotations of Python functions, arguments, and variables in a way
that can be used by various tools. It is an optional feature of Python
that has been introduced in Version~3.5 and enhanced in successive
versions. With this feature implemented a type checker, a tool
separate from the Python interpreter, can be used to statically
analyze code base to find bugs during early development stages and
spelunk code of large projects. Naturally, it can be used by IDEs to
interactively provide hints and suggestions to programmers and type
check the code as it is being developed. Annotations can and should be
placed near the code if possible. Observe that the Python interpreter
ignores these annotations. Annotations can also be placed in external
files (that have the extension \lstinline{pyi}), referred to as stubs,
for cases where the code base is inaccessible, such as the case with
bindings. Stub files contain the signatures of the annotated functions
with the body of the functions discarded. They also list the
attributes of annotated types. In our case they contain the annotated
interface of the \cgal{} bindings. Stub files have the same syntax as
regular Python modules. There is one feature of the typing module that
is different in stub files, namely, the \pyLstinline{@overload}
decorator. This decorator allows describing functions that support
multiple different combinations of argument types. Consider, for
example, the overloaded free C++ functions \cppLstinline{overlay()};
see Section~\ref{sec:code:overlay}. The annotation for the
corresponding Python function \pyLstinline{overlay()} in the stub file
is shown in Listing~\ref{lst:py:overloading}.

\begin{lstfloat}
\begin{lstlisting}[style=pythonFramedStyle,basicstyle=\normalsize,%
                   label={lst:py:overloading},%
                   caption={Type annotation of the overloaded overlay function.}]
@overload
def overlay(r:Arrangement_2, b:Arrangement_2, o:Arrangement_2) -> None: pass
@overload
def overlay(r:Arrangement_2, b:Arrangement_2, o:Arrangement_2, t:Arr_overlay_traits) -> None: pass
@overload
def overlay(r:Arrangement_2, b:Arrangement_2, o:Arrangement_2, t:Arr_overlay_function_traits) -> None: pass
\end{lstlisting}
\end{lstfloat}

\begin{mdframed}[hidealllines=true,backgroundcolor=red!20]
Annotating a Python type that wraps a C++ type, which is a model of
some concepts is guided by the concepts. Similar to the C++ code of
the bindings, we have introduced a tight coupling between concepts and
Python type annotations; see Section~\ref{sec:code:concepts}. In
particular, we have introduced a Python annotation class for each
concept that annotates all the requirements of the concept. The
refinement relations between concepts are not reflected using
inheritance between concept annotation classes.
% (This would invalidate the annotation code, as a concept \cppLstinline{A} may refine concepts \cppLstinline{B} and \cppLstinline{C} and both, \cppLstinline{B} and \cppLstinline{C}, may refine concept \cppLstinline{D}.)
Instead, every model annotation-class that annotates a certain type,
say \cppLstinline{T}, inherits from all the concept annotation-classes
that correspond to the concepts that \cppLstinline{T} models. For
example, assume that bindings are generated for the
\cppLstinline{Arrangement_2<GeometryGtaits_2, Dcel>} type instantiated
with a the geometry traits type that is an instance of the
\cppLstinline{Arr_segment_traits_2} class template. This type models
several concepts; see, e.g., Figure~\ref{fig:atc:central}. Partial
Python type annotation for the type \pyLstinline{Arr_segment_traits_2}
is shown in Listing~\ref{lst:py:aosSegmentTraits}.
\end{mdframed}

\begin{lstfloat}
  \begin{lstlisting}[style=pythonFramedStyle,basicstyle=\normalsize,%
                     backgroundcolor=\color{red!20},
                     label={lst:py:aosSegmentTraits},%
                     caption={Partial Python type annotation for the type \pyLstinline{Arr_segment_traits_2}.}]
class Arr_segment_traits_2(AosBasicTraits_2._AosBasicTraits_2,
                           AosXMonotoneTraits_2._AosXMonotoneTraits_2,
                           AosTraits_2._AosTraits_2,
                           AosLandmarkTraits_2._AosLandmarkTraits_2,
                           AosVerticalSideTraits_2._AosVerticalSideTraits_2,
                           AosVerticalSideTraits_2._AosHorizontalSideTraits_2) :
  class Point_2(AosBasicTraits_2._AosBasicTraits_2.Point_2):
    def x(self) -> Ker.FT: ...
    def y(self) -> Ker.FT: ...
  \end{lstlisting}
\end{lstfloat}

\begin{mdframed}[hidealllines=true,backgroundcolor=red!20]
The type annotation for the concept \cppLstinline{AosBasicTraits_2} contains the annotations for the default constructor and copy constructor of the type \cppLstinline{Point_2} defined by every model of the concept \cppLstinline{AosBasicTraits_2}; Listing~\ref{lst:py:aosBasicTraits}.
\end{mdframed}

\begin{lstfloat}
  \begin{lstlisting}[style=pythonFramedStyle,basicstyle=\normalsize,%
                     backgroundcolor=\color{red!20},
                     label={lst:py:aosBasicTraits},%
                     caption={Partial type annotation for the attributes that correspond to the concept \cppLstinline{AosBasicTraits_2}.}]
class _AosBasicTraits_2():
  class Point_2():
    def __init__(): ...
  class X_monotone_curve_2():
    def __init__(): ...
  class Equal_2():
    @overload
    def __call__(self, p: _AosBasicTraits_2.Point_2, q: _AosBasicTraits_2.Point_2) -> bool: ...
    @overload
    def __call__(self, p: _AosBasicTraits_2.X_monotone_curve_2, q: _AosBasicTraits_2.X_monotone_curve_2) -> bool: ...
  def equal_2_object(self) -> _AosBasicTraits_2.Equal_2: ...
  \end{lstlisting}
\end{lstfloat}
